% Book-owned openings and handoffs. The chapter sources remain independently
% buildable; these rails exist only in the Book. They are keyed by chapter
% prefix (see book/textbook.json), never by position, so reordering the
% book moves a chapter's rails with it. The handoff text, however, is written
% for the order the book is in: it names the chapter that comes next.
%
% An opening rail sits at the top of a chapter's first body page. Its heading
% is the chapter's question (from textbook.json, the same line the opener page
% set large); its body is the lede: the scene, in two or three sentences,
% before the chapter's own first section begins.

\newcommand{\pdbookrail}[2]{%
  \par\addvspace{12pt}%
  % A short opening/handoff is one paragraph, not a page-turn destination
  % for its last two lines. Keep the accent and its complete prose together.
  \noindent\begin{minipage}{\linewidth}
  % Keep question and handoff prose free of stretched interword spacing.
  \raggedright
  \noindent{\color{\pdcurrentchaptercolor}\rule{2.5em}{1.5pt}}\par
  \vspace{0.12cm}%
  \noindent{\small\scshape\bfseries\color{hhink}#1}\enspace
  {\small #2}\par
  \end{minipage}\par\addvspace{12pt}%
}

% ---- Part I: The Machine -------------------------------------------------

\newcommand{\pdchapteropeningswk}{%
  \pdbookrail{\pdchapterquestion}{%
    Alice's workers have two repairs based on revision~7. The first question is
    whether the second may still write after the first commits. This chapter
    locates the admission boundary, traces a stale request and a crash, and
    distinguishes a recorded conflict from a prevented effect.}%
}

\newcommand{\pdchapterhandoffswk}{%
  \pdbookrail{What \pdchapref{anchor}{The Anchor Protocol} receives}{%
    The repair now has a work-unit record, version checks and a stated recovery
    contract. These mechanisms distinguish a current proposal from a stale one;
    they do not determine who owns the repository. Next Bob must grant Alice
    precisely the authority needed for the job, including what her workers may
    inherit.}%
}

\newcommand{\pdchapteropeninganchor}{%
  \pdbookrail{\pdchapterquestion}{%
    Bob permits a checkout repair, not an unrelated migration or deployment.
    Alice delegates part of that work to Heron. This chapter follows the grant
    through signing, narrowing, presentation and revocation. The receiver must
    verify both the credential and its authority over the requested resource.}%
}

\newcommand{\pdchapterhandoffanchor}{%
  \pdbookrail{What \pdchapref{sealed}{The Sealed Harbor} receives}{%
    A verified grant can identify the permitted resource and action. It does
    not answer every disclosure question. Suppose diagnosis needs Derek's
    private checkout records and Erin's private model: neither owner has given
    Bob or Alice unrestricted access. The next chapter specifies how an answer
    may leave while their inputs remain subject to separate release policies.}%
}

\newcommand{\pdchapteropeningsealed}{%
  \pdbookrail{\pdchapterquestion}{%
    Derek supplies confidential records; Erin supplies the diagnostic model.
    They want a useful checkout diagnosis without releasing either input to
    the other party. We first write down the contract both owners sign and trace the permitted
    output, then state the channel, host and verifier assumptions needed to
    justify the confidentiality claims.}%
}

\newcommand{\pdchapterhandoffsealed}{%
  \pdbookrail{What \pdchapref{ls}{The Legible Swarm} receives}{%
    Alice may receive a permitted diagnosis without seeing the private evidence
    behind it. Other evidence is visible but too voluminous to read in full.
    These are different limits. The next chapter asks which observations a
    supervisor needs to distinguish a real repair from a deleted test or an
    overbroad patch, and how a summary should direct that inspection.}%
}

% ---- Part II: The Operator -----------------------------------------------

\newcommand{\pdchapteropeningls}{%
  \pdbookrail{\pdchapterquestion}{%
    Drake deletes the failing test. Wren fixes the defect but includes an
    unrelated migration. Heron contributes an acceptable repair. All three can
    produce an encouraging status line. This chapter asks what Alice must
    inspect, which records a successor needs, and how those demands compete
    for limited attention. Each bound uses an explicit observation model.}%
}

\newcommand{\pdchapterhandoffls}{%
  \pdbookrail{What \pdchapref{stp}{From Spawn to Person} receives}{%
    The supervisor can now ask which artifact supports a claim and which
    uncertainty remains. Tomorrow, however, Heron's process may be gone.
    To attribute its contribution and recover what it still owes, the next
    chapter separates an execution instance from a persistent participant and
    the principal responsible for it.}%
}

% ---- Part III: The Person ------------------------------------------------

\newcommand{\pdchapteropeningstp}{%
  \pdbookrail{\pdchapterquestion}{%
    Heron's process exits before all of its work is settled. A replacement can
    load the patch, but a new process label does not say which obligations or
    sanctions remain. This chapter defines the continuity needed to connect
    those records, then tests what reputation can and cannot mean under reset,
    migration and changing execution models.}%
}

\newcommand{\pdchapterhandoffstp}{%
  \pdbookrail{What \pdchapref{he}{The Harbor Economy} receives}{%
    A persistent participant links the repair to a responsible principal and
    retained outcomes. It does not make the work valuable or the reports true.
    The next chapter gives Bob and Alice an explicit purchase contract, with
    separate payment, collateral, evidence and settlement rules.}%
}

% ---- Part IV: The Market -------------------------------------------------

\newcommand{\pdchapteropeninghe}{%
  \pdbookrail{\pdchapterquestion}{%
    Bob buys the repair from Alice, who may also lease a runtime and use a
    metered skill. A labor price alone does not determine what each party
    receives after success, partial delivery or failure. This chapter follows
    those amounts through a worked ledger before generalizing the allocation,
    verification and incentive questions.}%
}

\newcommand{\pdchapterhandoffhe}{%
  \pdbookrail{What \pdchapref{bonded}{The Bonded Commons} receives}{%
    The contract names the parties, reserves and settlement rule. A disputed
    result still needs a procedure: who may submit evidence, who decides, and
    what remedy the decision authorizes. The next chapter keeps those questions
    separate from the immediate recovery of a worker that has gone silent.}%
}

\newcommand{\pdchapteropeningbonded}{%
  \pdbookrail{\pdchapterquestion}{%
    Heron's heartbeat stops while its contribution is being reviewed. Alice
    needs to recover the work; Bob may also contest whether the contract was
    satisfied. Silence alone decides neither question. This chapter develops
    recovery, adjudication and settlement as distinct procedures, then checks
    their stated accounting and incentive conditions.}%
}

\newcommand{\pdchapterhandoffbonded}{%
  \pdbookrail{What \pdchapref{fh}{The Federated Harbor} receives}{%
    The local institution now has separate rules for admission, recovery,
    evidence, adjudication and settlement. The repair must next run on a
    staging host governed by another operator. The final chapter asks which
    records travel, which authority must be granted there, and which
    disagreements the available observations can expose.}%
}

\newcommand{\pdchapteropeningfh}{%
  \pdbookrail{\pdchapterquestion}{%
    A staging operator receives the repair's signed receipts. They can help
    establish what an identified sender reported; they do not grant access to
    the staging machine. This chapter follows local authorization across the
    boundary, then asks when communication paths can witness inconsistent
    reports and when they leave a fault unobservable.}%
}

\newcommand{\pdchapterhandofffh}{%
  \pdbookrail{The conclusion of the book}{%
    The checkout job has acquired a sequence of inspectable contracts: what
    may change, who may authorize it, what may be disclosed, what evidence a
    reader needs, what survives replacement, how money is conserved, and who
    decides a dispute. Each contract preserves its own assumptions. Together
    they provide questions and mechanisms for designing accountable work;
    the book does not establish that every proposed component has been
    deployed or that their composition removes every failure.}%
}
